\section{Thought experiments}

\subsection{The scribe}
A scribe writes one digit every Planck instant --- nothing in the Universe changes faster.
Since the Big Bang he would have produced about \(10^{60}\) digits. Level \(b\) has
\(\bvalue\) of them. He has not finished even a negligible fraction; he will never reach the
end of \(b\) --- let alone \(c\). And \(b\) is the smallest of our giants. If even time does
not reach the end of \(b\), what does it mean to say that \(b\) has ``so many'' digits?

\subsection{The Library of Babel}
Borges's library holds every 410-page book --- conceivably vast, yet still finite. Beside
\(d\) it is laughably small. To shelve the digits of \(d\) you would need a library whose
catalogue could not be written in this Universe. Borges imagined the largest library a human
can bear; \(d\) demands one that cannot be borne.

\subsection{One who could see it}
Imagine a mind able to hold the crypto coma whole --- all its digits at once. Would the
crypto coma be, for it, a single number, the way seven is for us? Or does ``a single
number'' cease to mean anything at that scale? Perhaps to truly see a number is to make it
small. And is understanding not just another word for ``to shrink''?
